% Numbers from results-process.json (analyze_process.py).
Sixteen tasks in Repositories A and B (eight per repository): \emph{ticket-to-code} (which top-level functions did a
ticket's commits change that still exist, and which test files call them; truth from git history and import
resolution) and \emph{commit-to-trace} (a commit's ticket, that ticket's epic and specification; truth from the
ticket's own trace table). Arms plain, core and pack, five runs each (240 runs, all valid).

\paragraph{H6 not supported} Packs cost 60\% less than plain [48, 70] ($p<10^{-4}$), answering in one turn against
seven, but scored 0.080 lower [$-0.094$, $-0.058$], outside the non-inferiority margin. Plain search was nearly perfect
on these tasks (mean score 0.99). The full tool set was no cheaper ($-7\%$ [$-23$, 7]) and scored 0.12 lower.

\paragraph{Where it failed} On commit-to-trace, packs matched plain exactly (score 1.0) at 37\% less cost. On
ticket-to-code, the loss came from two of eight tasks. In one, the graph did not link the ticket to its commits; in
the other, the pack's test list missed most of the test files. Each time the agent accepted the pack in one turn and
did not search. This is the third exploratory mechanism, complete-looking lists, recurring in a pre-registered
setting. The ticket pack's coverage statement reports call sites it could not attribute, but not commits it failed
to link to the ticket, so a missed join looked like a complete answer. Coverage has to extend to every join a pack
makes.
